% Generated draft: requires mathematical and source review.
\documentclass[11pt]{article}
\usepackage[margin=25mm]{geometry}
\usepackage{fontspec}
\setmainfont{Latin Modern Roman}
\usepackage{amsmath,amssymb,mathtools,unicode-math}
\setmathfont{Latin Modern Math}
\usepackage{enumitem,needspace,xurl,hyperref}
\setlength{\parindent}{0pt}
\setlength{\parskip}{5pt}
\setlength{\emergencystretch}{3em}
\newcommand{\N}{\mathbb{N}}
\newcommand{\Z}{\mathbb{Z}}
\newcommand{\Q}{\mathbb{Q}}
\newcommand{\R}{\mathbb{R}}
\newcommand{\C}{\mathbb{C}}
\newcommand{\F}{\mathbb{F}}
\newcommand{\A}{\mathbb{A}}
\newcommand{\PP}{\mathbb{P}}
\newcommand{\E}{\mathbb{E}}
\newcommand{\eps}{\varepsilon}
\newcommand{\dd}{\,\mathrm{d}}
\begin{document}
\section*{UnsolvedMath ID: 10400054}
\textbf{Source problem code:} catalogue entry
\subsection{Definitions and mathematical statement}
Let $K\subset S^3$ be an oriented knot, and let $-K$ be the same knot with its orientation reversed. This is not a mirror-image operation. Let $\mathcal A(S^1)$ be the source's graded algebra of Jacobi diagrams on an oriented circle, with product given by connected sum along the circle, and let $\widehat{\mathcal A}(S^1)$ be its degree completion. Write $Z(K)$ for the Kontsevich invariant with the source's fixed diagram/framing normalization.

Prove that every oriented knot satisfies
\[
Z(K)=Z(-K)\quad\text{in }\widehat{\mathcal A}(S^1).
\]
The equality is in every degree, not just a fixed truncation. The source cross-refers to Conjecture 2.7 for an equivalent assertion, whose statement is not included and is not reconstructed.

The following text in the extract is background from Section 3.3: $Z(K)$ is group-like, so with its unit constant term the formal logarithm
\[
\log Z(K)=\sum_{r\geq1}\frac{(-1)^{r+1}}r(Z(K)-1)^r
\]
lies in the completed connected-diagram subspace, spanned by diagrams whose uni-trivalent part is connected. This structural fact is not an additional request to characterize all possible Kontsevich invariants.
\subsection{Short English statement}
Show that the Kontsevich invariant cannot distinguish any oriented knot from the same knot with reversed orientation, in every diagram degree.
\subsection{Sources}
\begin{itemize}
\item Ohtsuki, \emph{Problems on invariants of knots and 3-manifolds} (2002), Conjecture 3.4, PDF page 53, with its cross-reference to Conjecture 2.7 and following Section 3.3 background, \url{https://msp.org/gtm/2002/04/gtm-2002-04-024s.pdf}; accessed 2026-07-29, PDF-font-structure extraction.
\item Local registry \texttt{data/unsolvedmath/problems.json}, permanent UnsolvedMath ID $10400054$, source code catalogue entry, \emph{Conjecture 3.4 --- Orientation reversal of the Kontsevich invariant}. Dataset attribution: Ulam AI and the UnsolvedMath Contributors.
\end{itemize}
\end{document}
